\documentclass[oneside,12pt]{exam}          % ----- exam

%%%%%%%%%%%%%%%%%%%%%%%%%%%%
%%%%% Police et langue %%%%%
%%%%%%%%%%%%%%%%%%%%%%%%%%%%

\usepackage[x11names,table]{xcolor}                           % ----- Couleurs
\usepackage[utf8]{inputenc}                             % ----- UTF-8
\usepackage[french]{babel}                              % ----- French
\usepackage{amsmath,amssymb,mathrsfs}               % ----- Maths
\usepackage{fontawesome5,txfonts}						% ----- Caractères / Symboles
\usepackage{eurosym}                                    % ----- Currencies

\usepackage{helvet}										% ----- Police texte
\renewcommand{\familydefault}{\sfdefault}
\usepackage[T1]{fontenc}

\usepackage{mathpazo}									% ----- Police maths


%%%%%%%%%%%%%%%%
%%%%% TikZ %%%%%
%%%%%%%%%%%%%%%%

\usepackage{pgfplots}                                   % ----- Plots
    \pgfplotsset{compat=1.18}
\usepackage{tikz, tikz-3dplot, tkz-tab}                 % ----- TikZ, TikZ 3D, Tableaux

    \usetikzlibrary{math,arrows,arrows.meta}                            % ----- Maths, Vecteurs
    \usetikzlibrary{automata, positioning,patterns.meta}    % ----- Graphs
    \usetikzlibrary{decorations.fractals}                   % ----- Fractals
    \usetikzlibrary{tikzmark}
    \usetikzlibrary{fadings}

\tikzset{arrow style/.style = {\tkzTabDefaultWritingColor,-Stealth,> = \tkzTabDefaultArrowStyle,shorten > = \tkzTabDefaultSep,shorten < = \tkzTabDefaultSep}}		% ----- Modifier flèche tabVar

\tikzset{node style/.style = {inner sep = \tkzTabDefaultSep, outer sep = \tkzTabDefaultSep}}

%%%%%%%%%%%%%%%%%%%%
%%%%% Pratique %%%%%
%%%%%%%%%%%%%%%%%%%%

\usepackage{titlesec}							% ----- Titres et sections
\usepackage{calc}                                       % ----- Calc / Minipage
\usepackage{fancyhdr}                                   % ----- Header and footer
\usepackage{geometry}                                   % ----- Géométrie de la page
\usepackage{tcolorbox}                                  % ----- Boxes environment
    \tcbuselibrary{skins,listings,listingsutf8}             % ----- All libraries
\usepackage{varwidth}                                   % ----- Adjust box width
\usepackage{fourier-orns}                               % ----- Déco
\usepackage{witharrows}									% ----- Équations avec flèches
\usepackage{array}										% ----- Environnements array / tabular
\newcolumntype{P}[1]{>{\centering\arraybackslash}p{#1}}	% ----- Colonne centrée de taille ajustable
\usepackage{enumitem}									% ----- Customisation des \enumerate
\usepackage[c]{esvect}									% ----- Vecteurs
\usepackage{diagbox}									% ----- Découper des cases de \tabular
\usepackage{esvect}
\usepackage{graphicx}									% ----- Insérer des images			% ----- Packages utiles
%%%%%%%%%%%%%%%%%%%%%%%%
%%%%% Mise en page %%%%%
%%%%%%%%%%%%%%%%%%%%%%%%

\usepackage{parskip}									% ----- Parskip
    \setlength{\parskip}{\medskipamount}
    \newcommand{\restoreskip}{\setlength{\parskip}{\medskipamount}}
    
\setlength\parindent{0mm}                               % ----- No indent

\usepackage{setspace}                                   % ----- Interligne
	\setstretch{0.95}                           % Interligne 0.95
\usepackage{makecell}



%%%%%%%%%%%%%%%%%%%%%%%%%%%%%%%
%%%%% Format des sections %%%%%
%%%%%%%%%%%%%%%%%%%%%%%%%%%%%%%

\titlespacing{\section}{0pt}{5mm}{0pt}
\titlespacing{\subsection}{0pt}{2mm}{1mm}


\titleformat{\section}
    {\large\normalfont\bfseries}
    {Exercice \no \thesection\, - \,}
    {0pt}
    {\large\bfseries}

\renewcommand{\thesubsection}{\Alph{subsection}}
\titleformat{\subsection}
    {\normalfont\bfseries}
    {Partie \thesubsection\,: \,}
    {0pt}
    {\bfseries}





		% ----- Préambule feuilles de TD
%%%%%%%%%%%%%%%%%%%%%%%%%%%%%
%%%% Lettres double barre %%%
%%%%%%%%%%%%%%%%%%%%%%%%%%%%%

\newcommand\CC{%
        {\mathbb C}} % ----- Ensemble C

\newcommand\RR{%
        {\mathbb R}} % ----- Ensemble R

\newcommand\QQ{%
        {\mathbb Q}} % ----- Ensemble Q

\newcommand\ZZ{%
        {\mathbb Z}} % ----- Ensemble Z

\newcommand\NN{%
        {\mathbb N}} % ----- Ensemble N

\newcommand\DD{%
        {\mathbb D}} % ----- Ensemble D
        
\newcommand\UU{%
        {\mathbb U}} % ----- Ensemble U

\newcommand\PP{%
        {\mathbb P}} % ----- P probabiliste

\newcommand\EE{%
        {\mathbb E}} % ----- E espérance

\newcommand\VV{%
        {\mathbb V}} % ----- V variance



%%%%%%%%%%%%%%%%%%%%%%%%%%%%%%%%
%%%% Lettres calligraphiques %%%
%%%%%%%%%%%%%%%%%%%%%%%%%%%%%%%%

\newcommand\Ccal{%
        {\mathscr{C}}} % ----- C calligraphique

\newcommand\Dcal{%
        {\mathscr{D}}} % ----- D calligraphique

\newcommand\Pcal{%
        {\mathscr{P}}} % ----- P calligraphique

\newcommand\Rcal{%
        {\mathscr{R}}} % ----- R calligraphique

\newcommand\Fcal{%
        {\mathscr{F}}} % ----- F calligraphique

\newcommand\Scal{%
        {\mathscr{S}}} % ----- S calligraphique

\newcommand\Acal{%
        {\mathscr{A}}} % ----- A calligraphique

\newcommand\Vcal{%
        {\mathscr{V}}} % ----- V calligraphique
        
\newcommand\Tcal{%
        {\mathscr{T}}} % ----- V calligraphique



%%%%%%%%%%%%%%%%%%%%%%%%%%%%%%%%%%%%%
%%% Symboles et/ou calculs usuels %%%
%%%%%%%%%%%%%%%%%%%%%%%%%%%%%%%%%%%%%

%%%%% Analyse %%%%%

\newcommand\limx{%
        {\lim\limits_{x\to +\infty}}\,\,} % ----- Limite x en + infini

\newcommand\limn{%
        {\lim\limits_{n\to +\infty}}\,\,} % ----- Limite n en + infini

\newcommand\limite[1]{%
        {\lim\limits_{\substack{#1}}}\,\,} % ----- Limite X vers Y

\newcommand\nN{%
        {n\in\mathbb N}} % ----- n dans N

\newcommand\xR{%
        {x\in\mathbb R}} % ----- x dans R

\newcommand\zC{%
        {z\in\mathbb C}} % ----- z dans C

\newcommand\somme[2]{%
        {\displaystyle\sum_{#1}^{#2}}}    % ----- Somme

\newcommand\sommek[2]{%
        {\displaystyle\sum_{k=#1}^{#2}}}    % ----- Somme k

\newcommand\inte[2]{%
        {\displaystyle\int_{#1}^{#2}}}   % ----- Intégrale

\newcommand\dd{%
        {\,\,\mathrm{d}}}       % ----- d différentiel

\newcommand\floor[1]{%
        {\lfloor #1 \rfloor}}   % ----- Partie entière inférieure

\newcommand\ceil[1]{%
        {\lceil #1 \rceil}}     % ----- Partie entière supérieure

\newcommand\abs[1]{%
        {\left\lvert #1 \right\rvert}}     % ----- Valeur absolue

\newcommand\pgcd{%
        {\mathrm{pgcd}}}       % ----- pgcd

\newcommand\ppcm{%
        {\mathrm{ppcm}}}       % ----- ppcm



%%%%% Probabilités %%%%%

\newcommand\barre[1]{%
    {\overline{#1}}} % ----- A barre



%%%%% Géométrie %%%%%

\newcommand\vcoord[2]{			% ----- Coordonnées d'un vecteur (Nom + coord)
	\renewcommand{\arraystretch}{1.1} \arraycolsep=0.5\arraycolsep
	\vv{#1}\mathbin{=}\begin{pmatrix}#2\end{pmatrix}
} 

\newcommand\Oij{\left(O\,;\,\vv{i}\,;\,\vv{j}\right)} % ----- (O;i;j)

\def\Oijk{\left(O\,;\,\Vec{i}\,;\,\Vec{j}\,;\,\Vec{k}\right)} % ---- (O;i;j;k)

\newcommand\norme[1]{%
        {\left\lVert #1 \right\rVert}}     % ----- Norme

\newcommand\normev[1]{%
        {\left\lVert \vv{#1} \right\rVert}}     % ----- Norme d'un vecteur

\newcommand\scal[2]{%
    {\vv{#1}\cdot\vv{#2}}} % ----- Produit scalaire
    
\newcommand\smeq{{\,=\,}}	% ----- Symbole '=' avec petit espace



%%%%% Algèbre %%%%%

\newcommand\Reel{				% ----- Partie réelle
    {\mathfrak{Re}}} 

\newcommand\Imag{				% ----- Partie imaginaire
    {\mathfrak{Im}}} 

\newcommand\matrice[1]{{ 				% ----- Matrice
    \renewcommand{\arraystretch}{0.7}
    \arraycolsep=0.7\arraycolsep\ensuremath{\begin{pmatrix} #1 \end{pmatrix}}}}

\newcommand\Ouv{(O\,;\,\vv{u}\,;\,\vv{v})} 	% ----- Plan complexe (O;u;v)



%%%%%%%%%%%%%%%%%%%%
%%% Mise en page %%%
%%%%%%%%%%%%%%%%%%%%

%%%%% Flèche SSI longueur adaptative %%%%%

\makeatletter
\newcommand\xLongleftrightarrow[2][]%
  {\ext@arrow 0099{\Longleftrightarrowfill@}{#1}{#2}}
\def\Longleftrightarrowfill@
  {\arrowfill@\Leftarrow\Relbar\Rightarrow}
\makeatother



%%%%% Adaptation du stretch des environnements 'cases' %%%%%

\makeatletter
\renewcommand*\env@cases[1][1.2]{%
  \let\@ifnextchar\new@ifnextchar
  \left\lbrace
  \def\arraystretch{#1}%
  \array{@{}l@{\quad}l@{}}%
}
\makeatother



%%%%% Ajustement hauteur de \ldots %%%%%
\let\oldldots\ldots
\renewcommand{\ldots}{
\parbox[b][8pt]{20pt}{\oldldots}
}                 % ----- Macros mathématiques
%%%%%%%%%%%%%%%%%%%%%%%%%%%%%%%%
%%%%% ----- Packages ----- %%%%%
%%%%%%%%%%%%%%%%%%%%%%%%%%%%%%%%

\usepackage{tcolorbox}                                  % ----- Boxes environment
    \tcbuselibrary{skins,listings,listingsutf8,breakable}             % ----- All libraries
\usepackage[bold=0.7]{xfakebold}


%%%%%%%%%%%%%%%%%%%%%%%%%%%%%%%%%%
%%%%% ----- New colors ----- %%%%%
%%%%%%%%%%%%%%%%%%%%%%%%%%%%%%%%%%

\definecolor{persianblue}{HTML}{221177} % ----- New blues
\definecolor{lightblue}{HTML}{EEFFFF}
\definecolor{softblue}{HTML}{6655BB}

\definecolor{maroon}{HTML}{992222}      % ----- New reds
\definecolor{lightred}{HTML}{FAEEEE}

\definecolor{calgreen}{HTML}{005500}    % ----- New greens
\definecolor{lightgreen}{HTML}{EEFFEE}
\definecolor{softgreen}{HTML}{449944}

\definecolor{brickorange}{HTML}{CC7033} % ----- New oranges
\definecolor{lightorange}{HTML}{FFF2E6}

\definecolor{borderpink}{HTML}{aa00bb} % ----- New pinks
\definecolor{lightpink}{HTML}{ffeeff}

\definecolor{newred}{HTML}{CC0000}     % ----- Colors for geometry
\definecolor{newblue}{HTML}{0000CC}
\definecolor{newgreen}{HTML}{00AA00}
\definecolor{neworange}{HTML}{FFAA00}







%%%%%%%%%%%%%%%%%%%%%%%%%%%%%%%%%%%%%%%%
%%%%% ----- \begin{theoreme} ----- %%%%%
%%%%%%%%%%%%%%%%%%%%%%%%%%%%%%%%%%%%%%%%

\newtcolorbox{theoreme}[2][]{%
    enhanced , breakable , parbox=false ,
    arc=5pt,
    frame style={left color=persianblue, right color=softblue},
    colback=lightblue,
    fonttitle=\bfseries,
    title={\faBook~~ \underline{Théorème} : #2~~\faBook},
    attach boxed title to top left={xshift=5mm,yshift=-2mm},
    boxed title style={
        colback=persianblue,
        colframe=persianblue
    },
    #1
}



%%%%%%%%%%%%%%%%%%%%%%%%%%%%%
%%%%% \begin{propriete} %%%%%
%%%%%%%%%%%%%%%%%%%%%%%%%%%%%

\newtcolorbox{propriete}[2][]{%
    enhanced , breakable , parbox=false ,
    arc=5pt,
    frame style={left color=persianblue, right color=softblue},
    colback=lightblue,
    fonttitle=\bfseries,
    title={\faBook~~ \underline{Propriété} : #2~~\faBook},
    attach boxed title to top left={xshift=5mm,yshift=-2mm},
    boxed title style={
        colback=persianblue,
        colframe=persianblue
    },
    #1
}



%%%%%%%%%%%%%%%%%%%%%%%%%%%%%%%%%%%%%%%%%%
%%%%% ----- \begin{definition} ----- %%%%%
%%%%%%%%%%%%%%%%%%%%%%%%%%%%%%%%%%%%%%%%%%

\newtcolorbox{definition}[2][]{%
    enhanced , breakable , parbox=false ,
    arc=5pt,
    frame style={left color=calgreen, right color=softgreen},
    colback=lightgreen,
    fonttitle=\bfseries,
    title={\faUniversity~~ \underline{Définition} : #2~~\faUniversity},
    attach boxed title to top left={xshift=5mm,yshift=-2mm},
    boxed title style={
        colback=calgreen,
        colframe=calgreen
    },
    #1
}



%%%%%%%%%%%%%%%%%%%%%%%%%%%%%%%%%%%%%%
%%%%% ----- \begin{astuce} ----- %%%%%
%%%%%%%%%%%%%%%%%%%%%%%%%%%%%%%%%%%%%%

\newtcolorbox{astuce}[1][]{%
    enhanced , parbox=false ,
    arc=10pt,
    frame hidden,
    colback=lightred,
    fonttitle=\bfseries,
    borderline={1mm}{0mm}{maroon,dotted},
    title={\faLightbulb~~\underline{Astuce}~~\faLightbulb},
    attach boxed title to top left={xshift=5mm,yshift=-2mm},
    boxed title style={
        colback=maroon,
        colframe=maroon
    },
    #1
}



%%%%%%%%%%%%%%%%%%%%%%%%%%%%%%%%%%%%%%%%
%%%%% ----- \begin{remarque} ----- %%%%%
%%%%%%%%%%%%%%%%%%%%%%%%%%%%%%%%%%%%%%%%

\newtcolorbox{remarque}[1][]{%
    enhanced , parbox=false ,
    arc=10pt,
    frame hidden,
    colback=lightred,
    fonttitle=\bfseries,
    borderline={1mm}{0mm}{maroon,dotted},
    title={\faExclamation~~\underline{Remarque}~~\faExclamation},
    attach boxed title to top left={xshift=5mm,yshift=-2mm},
    boxed title style={
        colback=maroon,
        colframe=maroon
    },
    #1
}



%%%%%%%%%%%%%%%%%%%%%%%%%%%%%%%%%%%%%%%
%%%%% ----- \begin{methode} ----- %%%%%
%%%%%%%%%%%%%%%%%%%%%%%%%%%%%%%%%%%%%%%

\newtcolorbox{methode}[2][]{%
    enhanced , parbox=false ,
    arc=10pt,
    frame hidden,
    colback=lightred,
    fonttitle=\bfseries,
    borderline={1mm}{0mm}{maroon,dotted},
    title={\faCogs~~\underline{Méthode} : #2~~\faCogs},
    attach boxed title to top left={xshift=5mm,yshift=-2mm},
    boxed title style={
        colback=maroon,
        colframe=maroon
    },
    #1
}



%%%%%%%%%%%%%%%%%%%%%%%%%%%%%%%%%%%%%%%
%%%%% ----- \begin{exemple} ----- %%%%%
%%%%%%%%%%%%%%%%%%%%%%%%%%%%%%%%%%%%%%%

\newcommand{\exempledecoration}{%
    \draw[color=brickorange , line width = 2pt , fill=white]
        (frame.south west) ++(0,1pt) -- ++(3cm,0) ++(0,1.5mm) rectangle ++ (3mm,-3mm) ;
}

\newtcolorbox{exemplebox}[1][]{%
	enhanced , breakable , parbox=false ,
    arc=0pt,
    frame hidden,
    borderline west={2pt}{0pt}{brickorange},
    interior style={
        top color=lightorange,
        bottom color=lightorange!15!white
    },
    fonttitle=\bfseries,
    overlay unbroken={\exempledecoration},
    overlay last={\exempledecoration},
    #1
}

\newenvironment{exemple}[1]{%
    \begin{exemplebox}
        \textbf{\color{brickorange}\underline{Exemple} : #1}\\[5pt]
}{%
    \end{exemplebox}
}




%%%%%%%%%%%%%%%%%%%%%%%%%%%%%%%%%%%%%%%%%%%%%
%%%%% ----- \begin{demonstration} ----- %%%%%
%%%%%%%%%%%%%%%%%%%%%%%%%%%%%%%%%%%%%%%%%%%%%

\newcommand{\demodecoration}{%
    \draw[color=borderpink , line width = 2pt , fill=white]
        (frame.south west) ++(0,1pt) -- ++(3cm,0) ++(0,1.5mm) rectangle ++ (3mm,-3mm) ;
}

\newtcolorbox{demobox}[1][]{%
	enhanced , breakable , parbox=false ,
    arc=0pt,
    frame hidden,
    borderline west={2pt}{0pt}{borderpink},
    interior style={
        top color=lightpink,
        bottom color=lightpink!15!white
    },
    fonttitle=\bfseries,
    overlay unbroken={\demodecoration},
    overlay last={\demodecoration},
    #1
}

\newenvironment{demonstration}{%
    \begin{demobox}
        \textbf{\color{borderpink}\underline{Démonstration}}\\[5pt]
}{%
    \end{demobox}
}



%%%%%%%%%%%%%%%%%%%%%%%%%%%%%%%%%
%%%%% ----- \exercice ----- %%%%%
%%%%%%%%%%%%%%%%%%%%%%%%%%%%%%%%%

\newtcolorbox[auto counter , number within=section , number freestyle={\noexpand\arabic{\tcbcounter}} ]{exercicebox}[2][]{% ----- Exercice box
enhanced,
arc=1pt,
colframe=black!80,
colback=black!10,
fonttitle=\bfseries,
hbox,
title={Exercice n° \thetcbcounter \, : #2 },
attach boxed title to top*,
boxed title style = {colback = black!80},
#1}


\newcommand\exercice[2]{% ----- Exercice command
    \begin{exercicebox}[]{#1}
        \begin{varwidth}{\linewidth-30pt} 
         #2
        \end{varwidth}
    \end{exercicebox}
}



%%%%%%%%%%%%%%%%%%%%%%%%%%%%%%%%
%%%%% ----- \calculs ----- %%%%%
%%%%%%%%%%%%%%%%%%%%%%%%%%%%%%%%

\newtcolorbox{calculbox}[1][]{% ----- Calculs box
enhanced,
arc=1pt,
colbacktitle=black!25!white,
coltitle=black,
colback=white,
colframe=black!25!white,
hbox,
title={Aides de calcul},
attach boxed title to top*,
#1}


\newcommand\calculs[1]{% ----- Calculs command
    \begin{calculbox}
        \restoreskip
        #1
    \end{calculbox}
}






%%%%%%%%%%%%%%%%%%%%%%%%
%%%%% Mise en page %%%%%
%%%%%%%%%%%%%%%%%%%%%%%%

%%%%% Flèche SSI longueur adaptative %%%%%

\makeatletter
\newcommand\xLongleftrightarrow[2][]%
  {\ext@arrow 0099{\Longleftrightarrowfill@}{#1}{#2}}
\def\Longleftrightarrowfill@
  {\arrowfill@\Leftarrow\Relbar\Rightarrow}
\makeatother



%%%%% Adaptation du stretch des environnements 'cases' %%%%%

\makeatletter
\renewcommand*\env@cases[1][1.2]{%
  \let\@ifnextchar\new@ifnextchar
  \left\lbrace
  \def\arraystretch{#1}%
  \array{@{}l@{\quad}l@{}}%
}
\makeatother



%%%%% Ajustement hauteur de \ldots %%%%%
\let\oldldots\ldots
\renewcommand{\ldots}{
\parbox[b][8pt]{20pt}{\oldldots}
}

 		   	% ----- Macros de mise en page

\usepackage{afterpage}

\geometry{left=30pt,right=30pt, top=100pt, bottom=55pt}


\begin{document}

\afterpage{\newgeometry{left=30pt,right=30pt, top=55pt, bottom=55pt}}


\firstpageheader{Prénom : ..................................................\\[10mm] Nom : ....................................................... }{}{Classe : ......................... \\[10mm] Date : .........................}

\footer{}{\footnotesize Lycée polyvalent Irène et Frédéric Joliot Curie}{}


\hrulefill\raisebox{-1.9pt}{\quad\decofourleft~\decosix~\decofourright \quad}\hrulefill

\begin{center}
	\textbf{\large Automatismes \no 1 (Évalué)}
\end{center}


%%%%%%%%%%%%%%%%%%%%%%%%%%%%%%%%%%%%%%%%%%%%%%%%%%%%%%%%%%%%%%%%%%%%%%%%%%%%%%%%%%%%%%%%%%%%%%%%%%%%%%%%%%%%%
%%%%%%%%%%%%%%%%%%%%%%%%%%%%%%%%%%%%%%%%%%%%%%%%%%%%%%%%%%%%%%%%%%%%%%%%%%%%%%%%%%%%%%%%%%%%%%%%%%%%%%%%%%%%%
%%%%%%%%%%%%%%%%%%%%%%%%%%%%%%%%%%%%%%%%%%%%%%%%%%%%%%%%%%%%%%%%%%%%%%%%%%%%%%%%%%%%%%%%%%%%%%%%%%%%%%%%%%%%%


Pour chacune des questions suivantes, une et une seule des réponses proposées est exacte. Pour répondre, entourez sur la feuille la réponse choisie. Aucune justification n'est demandée.


\begin{enumerate}[label=\underline{Question \no\arabic*} :  , leftmargin=* , itemsep=5mm]

%%%%% QUESTION 1 %%%%%	
	\item Combien de solutions réelles admet l'équation $x^2 - 1\,234 = 0$ ?
	
		\hspace{\linewidth-\textwidth}
		\renewcommand{\arraystretch}{1.5}
		\begin{tabular}{| P{0.22\textwidth} | P{0.22\textwidth} | P{0.22\textwidth} | P{0.22\textwidth} |} \hline
			A & B & C & D \\ \hline
			Aucune & 1 & 2 & Autre \\ \hline
		\end{tabular}



%%%%% QUESTION 2 %%%%%
	\item
	\begin{minipage}{0.45\textwidth} \restoreskip
		On a représenté sur le graphique ci-contre les différentes pratiques des adhérents d'un club de sport.
		
		Parmi les adhérents, quelle proportion pratique le chess-boxing ?
	\end{minipage}
	\begin{minipage}{0.35\textwidth}
		\begin{tikzpicture}
			\pie[after number = , text = inside , sum = 300 , radius = 3]
			{90/{Roller\\Derby}, 120/{Chess\\Boxing}, 60/{Course\\en sac}, 30/{Aqua\\Poney}}
		\end{tikzpicture}
	\end{minipage}
	
	
		\hspace{\linewidth-\textwidth}
		\renewcommand{\arraystretch}{1.5}
		\begin{tabular}{| P{0.22\textwidth} | P{0.22\textwidth} | P{0.22\textwidth} | P{0.22\textwidth} |} \hline
			A & B & C & D \\ \hline
			$30\%$& $35\%$ & $40\%$ & $45\%$ \\ \hline
		\end{tabular}
		
		

%%%%% QUESTION 3 %%%%%		
	\item Quelle est l'unique solution réelle de l'équation $26x - 15 = 12x + 34$ ?
	
		\hspace{\linewidth-\textwidth}
		\renewcommand{\arraystretch}{1.75}
		\begin{tabular}{| P{0.22\textwidth} | P{0.22\textwidth} | P{0.22\textwidth} | P{0.22\textwidth} |} \hline
			A & B & C & D \\ \hline
			$\dfrac{2}{7}$ & $3,5$ & $\dfrac{-7}{2}$ & $-5,3$ \\ \hline
		\end{tabular}
		
		

%%%%% QUESTION 4 %%%%%
	\item Amir a des bonbons. Il en donne $30\%$ à sa petite sœur et mange $40\%$ du reste. Quelle proportion a-t-il mangé ?
	
		\hspace{\linewidth-\textwidth}
		\renewcommand{\arraystretch}{1.5}
		\begin{tabular}{| P{0.22\textwidth} | P{0.22\textwidth} | P{0.22\textwidth} | P{0.22\textwidth} |} \hline
			A & B & C & D \\ \hline
			$14\%$ & $22\%$ & $28\%$ & $34\%$ \\ \hline
		\end{tabular}



%%%%%%%%%%%%%%%%%%%%%%%%%%%%%%%%%%%%%%%%%%%%%%%%%
\newpage
\end{enumerate}
	\newgeometry{left=30pt,right=30pt, top=30pt, bottom=55pt}
\begin{enumerate}[resume*]
%%%%%%%%%%%%%%%%%%%%%%%%%%%%%%%%%%%%%%%%%%%%%%%%%

		

%%%%% QUESTION 5 %%%%%
	\item Un article soldé à $-40\%$ coûte $150$\euro. Quel était son prix initial ?
	
		\hspace{\linewidth-\textwidth}
		\renewcommand{\arraystretch}{1.5}
		\begin{tabular}{| P{0.22\textwidth} | P{0.22\textwidth} | P{0.22\textwidth} | P{0.22\textwidth} |} \hline
			A & B & C & D \\ \hline
			$190$\euro & $200$\euro & $230$\euro & $250$\euro \\ \hline
		\end{tabular}
		
		

%%%%% QUESTION 6 %%%%%
	\item Une quantité $Q$ augmente de $50\%$, puis diminue de $50\%$. \\
	On peut conclure que la quantité $Q$ :
	
		\hspace{\linewidth-\textwidth}
		\renewcommand{\arraystretch}{1.5}
		\begin{tabular}{| P{0.22\textwidth} | P{0.22\textwidth} | P{0.22\textwidth} | P{0.22\textwidth} |} \hline
			A & B & C & D \\ \hline
			A augmenté de $25\%$ & N'a pas évolué & A diminué de $25\%$ & Autre \\ \hline
		\end{tabular}	



%%%%% QUESTION 7 %%%%%
	\item Lors d'une sortie au musée, l'intendance du lycée paie les entrée pour une classe de 30 élèves accompagnée par 3 enseignants. Chaque entrée coûte $15$\euro\, et les élèves profitent d'une réduction de $20\%$ \\
	Quel est le prix de la visite au musée ?
	
		\hspace{\linewidth-\textwidth}
		\renewcommand{\arraystretch}{1.5}
		\begin{tabular}{| P{0.22\textwidth} | P{0.22\textwidth} | P{0.22\textwidth} | P{0.22\textwidth} |} \hline
			A & B & C & D \\ \hline
			$315$\euro & $360$\euro & $405$\euro & $440$\euro \\ \hline
		\end{tabular}



%%%%% QUESTION 8 %%%%%
	\item Une certaine quantité $Q$ augmente deux fois de $x$\%. À l'issue de ces deux évolutions, la quantité $Q$ a subi une hausse de $96$\%. \\
	Quelle est la valeur de $x$ ?
	
		\hspace{\linewidth-\textwidth}
		\renewcommand{\arraystretch}{1.5}
		\begin{tabular}{| P{0.22\textwidth} | P{0.22\textwidth} | P{0.22\textwidth} | P{0.22\textwidth} |} \hline
			A & B & C & D \\ \hline
			$48$ & $38$ & $40$ & $50$ \\ \hline
		\end{tabular}	
		
\end{enumerate}



\newpage

\begin{center}
	\textbf{\underline{\Large Correction}}
\end{center}


\begin{enumerate}[itemsep = 2mm]
	\item
	\begin{minipage}{0.5\textwidth} \restoreskip
		$x^2 - 1\,234 = 0 \Longleftrightarrow x^2 = 1\,234$
		
		On se demande donc s'il existe des nombres dont le carré vaut $1\,234$
		
		Il en existe \textbf{2} comme l'illustre le schéma ci-contre :
	\end{minipage} \hspace{15mm}
	\begin{minipage}{0.5\textwidth}
	    \begin{tikzpicture}[scale=0.65] 
	        \begin{axis}[xmin = -40, xmax = 40, ymin = -0, ymax = 1500,
	            grid = none, ticks = none ,
	            axis lines = middle, enlargelimits = { abs = 0.2 },
	            clip = false
	            ]
	        
	        \addplot[color=newblue , samples = 100 , domain = -38:38 , line width = 2pt]
	        {x^2} ;
	        
	        \draw[color=newred , dashed] (-35,0) node[below] {$-\sqrt{1234}$} -- (-35,1234) -- (0,1234) node[above right] {$1234$} -- (35,1234) -- (35,0) node[below] {$\sqrt{1234}$} ;
	        
	        \draw[color=newblue] (45,1300) node {$y=x^2$} ;

	        
	        \end{axis}
	    \end{tikzpicture}  
	\end{minipage}
	
	\item
	Le nombre total d'adhérents est : $120 + 90 + 30 + 60 = 300$ et il y a $120$ pratiquants du chess-boxing.
	
	Or $\dfrac{120}{300} = \dfrac{3 \times 40}{3 \times 100} = \dfrac{40}{100}$ \quad , donc $40\%$ des adhérents pratiquent le chess-boxing.
	
	\item
	$\begin{WithArrows}[format = cc]
    & 26x - 15 = 12x + 34 \Arrow{On ajoute 15 de chaque côté} \\
    \Leftrightarrow  & 26x = 12x + 49 \Arrow{On soustrait $12x$ de chaque côté} \\
    \Leftrightarrow  & 14x = 49 \Arrow{On divise par $14$ de chaque côté} \\
    \Leftrightarrow  & x = \dfrac{49}{14}  = \dfrac{7}{2} = 3,5 
    \end{WithArrows}$

	\item
	Après avoir donné $30\%$ de ses bonbons, il en reste $70\%$ à Amir. Il a donc mangé $40\%$ de ces $70\%$
	
	$\dfrac{40}{100} \times \dfrac{70}{100} = \dfrac{4 \times 7}{10 \times 10} = \dfrac{28}{100}$ \quad , Amir a donc mangé $28\%$ de ses bonbons.
	
	\item On note $P$ le prix initial de l'article. On sait que ce prix $P$ a subi une baisse de $40\%$, or diminuer de $40\%$ revient à multiplier par $\left(1 - \dfrac{40}{100}\right) = \dfrac{60}{100}$ \\
	On sait donc que $P \times \dfrac{60}{100} = 150$ , alors $P = 150 \times \dfrac{100}{60} = \dfrac{30 \times 5 \times 100}{30 \times 2} = \dfrac{5 \times 50 \times 2}{2} = 250$ 
	
    \item
    Augmenter de $50\%$ revient à multiplier par $\left(1 + \dfrac{50}{100}\right)$ \\
    Diminuer de $50\%$ revient à multiplier par $\left(1 - \dfrac{50}{100}\right)$
    
    $\left(1 + \dfrac{50}{100}\right) \times \left(1 - \dfrac{50}{100}\right) = 1,5 \times 0,5 = 0,75$
    
    La quantité $Q$ a donc diminué de $25\%$
    
    \item
    Les entrées au plein tarif pour les élèves coûtent $30 \times 15$\euro $\,= 450$\euro \\    
    On applique la réduction de $20\%$ : $450 \times \left(1 - \dfrac{20}{100}\right) = 450 \times 0,8 = 360$ \\    
    Les enseignants paient leur entrée au plein tarif, donc il faut rajouter $3 \times 45$\euro $\,=45$\euro
    
    Le prix de la sortie est donc de $360$\euro $\,+15$\euro $\,=405$\euro
    
    \item
	Augmenter de $40\%$ revient à multiplier par $\left(1 + \dfrac{40}{100}\right)$ \\
	Or \quad $\left(1 + \dfrac{40}{100}\right) \times \left(1 + \dfrac{40}{100}\right) = 1,4 \times 1,4 = 1,96$
    
\end{enumerate}




\end{document}

		